% Directly inputtable paper section. Author: ChatGPT.
\section{Generated five-term relations and the finite-boundary bridge}
\label{sec:steinberg-five-term-boundary-bridge}

\noindent\textbf{Author: ChatGPT.}
The valuation-contact construction of
Section~\ref{sec:steinberg-integer-finite-chain} provides a concrete
finite-support exterior surface for every positive reduced rational cell and
an analytic estimate for every finite chain satisfying an exact surface
boundary.  What remained was to derive that boundary from the relation
generated by five-term moves.  We supply that derivation here.  We first work
in a free divisor-symbol module, then pass to the smaller relation generated
only by positively realized rational moves.  The argument is the divisor
shadow of the standard pre-Bloch complex
\cite{DupontSah1982,Suslin1991,GoetteZickert2007}; no $K$-theoretic or
$abc$ estimate is imported from that theory.

\subsection{The free symbol chain and its boundary}

Let
\[
 \mathcal D=\mathbb N\to_0\mathbb Z,
 \qquad
 \mathcal E=(\mathbb N\times\mathbb N)\to_0\mathbb Z,
\]
and retain the ordered-coordinate exterior product
\[
 (A\wedge B)_{p,q}=A_pB_q-A_qB_p.
\]
Write $\mathcal S=\mathcal D\times\mathcal D$ for the set of divisor
symbols and let
\[
 \mathcal C_0=\mathbb Z^{(\mathcal S)}
\]
be the free integral module on those symbols.  Its basis element associated
with $(A,B)$ is denoted by $[A,B]$.  Define a linear map
\begin{equation}
 \partial:\mathcal C_0\longrightarrow\mathcal E,
 \qquad \partial[A,B]=A\wedge B.
 \label{eq:sftbb-boundary}
\end{equation}

For $X,U,Y,V,Z\in\mathcal D$, define the formal five-term chain
\begin{align}
 \mathfrak r(X,U,Y,V,Z)={}&[X,U]-[Y,V]+[Y-X,Z-X]\notag\\
 &-[Y+U-X-V,Z-X-V]+[U-V,Z-V].
 \label{eq:sftbb-five-chain}
\end{align}

\begin{theorem}[Five-term generators are boundary cycles]
\label{thm:sftbb-generator-cycle}
For all $X,U,Y,V,Z\in\mathcal D$,
\[
 \partial\mathfrak r(X,U,Y,V,Z)=0.
\]
\end{theorem}

\begin{proof}
By linearity and~\eqref{eq:sftbb-boundary}, the claimed boundary is
\begin{align*}
 X\wedge U-Y\wedge V\\
 &+(Y-X)\wedge(Z-X)\\
 &-(Y+U-X-V)\wedge(Z-X-V)\\
 +(U-V)\wedge(Z-V).
\end{align*}
Expand each wedge bilinearly.  Every elementary wedge between two members of
$\{X,U,Y,V,Z\}$ occurs with opposite coefficients and cancels.  This is
also the finite-support lift of the earlier coefficientwise five-term
identity.
\end{proof}

Let
\begin{equation}
 \mathcal R_5=\operatorname{span}_{\mathbb Z}
 \{\mathfrak r(X,U,Y,V,Z):X,U,Y,V,Z\in\mathcal D\}
 \subseteq\mathcal C_0.
 \label{eq:sftbb-relation}
\end{equation}

\begin{theorem}[Generated relation-to-boundary bridge]
\label{thm:sftbb-span-kernel}
One has
\[
 \mathcal R_5\subseteq\ker\partial.
\]
Consequently, if
\[
 c\sim_5 d
 \quad\Longleftrightarrow\quad c-d\in\mathcal R_5,
\]
then $\partial c=\partial d$.
\end{theorem}

\begin{proof}
The kernel of $\partial$ is an integral submodule.  By
Theorem~\ref{thm:sftbb-generator-cycle}, it contains every generator of the
span in~\eqref{eq:sftbb-relation}; hence it contains the span.  If
$c-d\in\mathcal R_5$, then linearity gives
$\partial c-\partial d=\partial(c-d)=0$.
\end{proof}

This proof includes arbitrary finite sums, negatives, and integer
multiples.  In particular, the exact boundary is a consequence of a
generation certificate and is not included among its premises.

\subsection{The positive rational subrelation}

For a positive reduced cell
\[
 P=(a,b,c),\qquad a+b=c,\qquad\gcd(a,b)=1,
\]
define its divisor symbol by
\begin{equation}
 \sigma(P)=\bigl(d(a)-d(c),d(b)-d(c)\bigr).
 \label{eq:sftbb-cell-symbol}
\end{equation}
Then
\begin{equation}
 \partial[\sigma(P)]
 =(d(a)-d(c))\wedge(d(b)-d(c))=\Omega(P).
 \label{eq:sftbb-cell-boundary}
\end{equation}

A positive realization of a five-term generator consists of positive
reduced cells $P_1,\ldots,P_5$ and divisor data $X,U,Y,V,Z$ for which
\begin{equation}
\begin{array}{c|c}
P_1 &(X,U)\\
P_2 &(Y,V)\\
P_3 &(Y-X,Z-X)\\
P_4 &(Y+U-X-V,Z-X-V)\\
P_5 &(U-V,Z-V)
\end{array}
\label{eq:sftbb-positive-realization}
\end{equation}
are the five individual identities for $\sigma(P_i)$.  These local symbol
identities contain no signed boundary assertion.  Let $\mathcal R_5^+$ be
the integral span of the signed chains
\begin{equation}
 [\sigma(P_1)]-[\sigma(P_2)]+[\sigma(P_3)]
 -[\sigma(P_4)]+[\sigma(P_5)]
 \label{eq:sftbb-positive-generator}
\end{equation}
over all positive realizations.

\begin{theorem}[Positive generated boundary]
\label{thm:sftbb-positive-span-kernel}
There are inclusions
\[
 \mathcal R_5^+\subseteq\mathcal R_5\subseteq\ker\partial.
\]
Thus every equivalence generated by actual positive five-term moves
preserves the finite-support surface.
\end{theorem}

\begin{proof}
The five identities in~\eqref{eq:sftbb-positive-realization} identify each
chain~\eqref{eq:sftbb-positive-generator} with the algebraic generator
\eqref{eq:sftbb-five-chain}.  Therefore each generator of
$\mathcal R_5^+$ belongs to $\mathcal R_5$, and the first inclusion follows
by taking spans.  The second inclusion is
Theorem~\ref{thm:sftbb-span-kernel}.
\end{proof}

Suppose now that a target cell $P_0$ and cells $P_j$ with coefficients
$n_j\in\mathbb Z$ satisfy the generation statement
\begin{equation}
 [\sigma(P_0)]-\sum_jn_j[\sigma(P_j)]\in\mathcal R_5^+.
 \label{eq:sftbb-generated-filling}
\end{equation}
Theorem~\ref{thm:sftbb-positive-span-kernel} and
\eqref{eq:sftbb-cell-boundary} derive
\begin{equation}
 \Omega(P_0)=\sum_jn_j\Omega(P_j).
 \label{eq:sftbb-derived-boundary}
\end{equation}
Applying the exact finite-chain theorem from the preceding section gives
\begin{equation}
 \Phi(P_0)\leq\frac12\left(
   \sum_j|n_j|\bigl(M(P_j)+V(P_j)\bigr)
   +\sum_j|n_j|R(P_j)\right).
 \label{eq:sftbb-generated-cost}
\end{equation}
Unlike the earlier conditional structure,
\eqref{eq:sftbb-generated-filling} stores only membership in the relation
generated by actual moves.  Equation~\eqref{eq:sftbb-derived-boundary} is a
theorem.

\subsection{An infinite arithmetic family of positive moves}

Let $0<b<a<c$ and impose the five reducedness conditions
\begin{equation}
\begin{gathered}
 \gcd(a,c-a)=\gcd(b,c-b)=\gcd(b,a-b)=1,\\
 \gcd\bigl(b(c-a),c(a-b)\bigr)=1,\\
 \qquad \gcd(c-a,a-b)=1.
\end{gathered}
\label{eq:sftbb-five-coprimalities}
\end{equation}
Consider the five cells
\begin{equation}
\begin{array}{c|c|c|c}
 &\text{left}&\text{middle}&\text{total}\\ \hline
P_1&a&c-a&c\\
P_2&b&c-b&c\\
P_3&b&a-b&a\\
P_4&b(c-a)&c(a-b)&a(c-b)\\
P_5&c-a&a-b&c-b.
\end{array}
\label{eq:sftbb-common-denominator-cells}
\end{equation}
Every entry is positive and the hypotheses
\eqref{eq:sftbb-five-coprimalities} make every row reduced.  The fourth row
is a cell because
\[
 b(c-a)+c(a-b)=a(c-b).
\]

\begin{proposition}[Common-denominator realization]
\label{prop:sftbb-common-denominator}
The cells~\eqref{eq:sftbb-common-denominator-cells} positively realize the
five-term move with
\[
 x=\frac ac,\qquad y=\frac bc,
\]
and arguments
\begin{equation}
 \frac ac,\quad\frac bc,\quad\frac ba,\quad
 \frac{b(c-a)}{a(c-b)},\quad\frac{c-a}{c-b}.
 \label{eq:sftbb-five-rationals}
\end{equation}
\end{proposition}

\begin{proof}
Put
\[
\begin{aligned}
X&=d(a)-d(c),&U&=d(c-a)-d(c),\\
Y&=d(b)-d(c),&V&=d(c-b)-d(c),\\
Z&=d(a-b)-d(c).
\end{aligned}
\]
The symbols of the first two cells are $(X,U)$ and $(Y,V)$ by definition.
Direct subtraction gives the third and fifth rows of
\eqref{eq:sftbb-positive-realization}.  For the fourth row,
multiplicativity of divisors gives
\begin{align*}
d\bigl(b(c-a)\bigr)-d\bigl(a(c-b)\bigr)&=Y+U-X-V,\\
d\bigl(c(a-b)\bigr)-d\bigl(a(c-b)\bigr)&=Z-X-V.
\end{align*}
Thus all five symbol identities hold, and
Theorem~\ref{thm:sftbb-positive-span-kernel} supplies the surface relation.
\end{proof}

For every odd $c\geq3$, the choice $a=2$, $b=1$ satisfies
\eqref{eq:sftbb-five-coprimalities}.  Indeed, the only nontrivial checks
reduce to
\[
 \gcd(2,c-2)=1,
 \qquad
 \gcd(c-2,c)=\gcd(c-2,2)=1.
\]
This gives an infinite parameterized family of actual positive generators.

\subsection{A full-premise counterexample to literal chain cancellation}

The generated quotient is essential.  Define the augmentation
\[
 \varepsilon:\mathcal C_0\longrightarrow\mathbb Z,
 \qquad
 \varepsilon\left(\sum_sn_s[s]\right)=\sum_sn_s.
\]
For every input one has
\begin{equation}
 \varepsilon\bigl(\mathfrak r(X,U,Y,V,Z)\bigr)
 =1-1+1-1+1=1.
 \label{eq:sftbb-augmentation}
\end{equation}

\begin{proposition}[A generator is never the zero free chain]
\label{prop:sftbb-generator-nonzero}
For every $X,U,Y,V,Z\in\mathcal D$,
\[
 \mathfrak r(X,U,Y,V,Z)\neq0.
\]
In particular, the boundary map $\partial$ is not injective.
\end{proposition}

\begin{proof}
If the free chain were zero, its augmentation would be zero, contradicting
\eqref{eq:sftbb-augmentation}.  Its boundary is nevertheless zero by
Theorem~\ref{thm:sftbb-generator-cycle}, so $\partial$ has a nonzero kernel.
\end{proof}

This refutes the exact shortcut asserting that the five signed basis chains
cancel before quotienting.  It does not refute the five-term route.  The
correct statement is the kernel inclusion of
Theorem~\ref{thm:sftbb-span-kernel}.

The failure persists with all positivity and reducedness premises.  Taking
\[
 x=\frac12,\qquad y=\frac16
\]
gives the five distinct arguments
\[
 \frac12,\quad\frac16,\quad\frac13,\quad\frac15,
 \quad\frac35
\]
and reduced cell triples
\[
 (1,1,2),\quad(1,5,6),\quad(1,2,3),\quad(1,4,5),
 \quad(3,2,5).
\]
Their signed surface boundary is zero, whereas their signed free chain is
nonzero; its coefficient at the $1/2$ symbol is one.

\subsection{Formal boundary and remaining analytic gate}

The Lean companion module is
\begin{center}
\texttt{SteinbergFiveTermBoundaryBridge20260903}.
\end{center}
It formalizes Theorems~\ref{thm:sftbb-generator-cycle},
\ref{thm:sftbb-span-kernel}, and
\ref{thm:sftbb-positive-span-kernel}, the generated filling implication
\eqref{eq:sftbb-generated-cost},
Proposition~\ref{prop:sftbb-common-denominator}, the odd-denominator family,
and Proposition~\ref{prop:sftbb-generator-nonzero}.  It also formalizes the
five-distinct positive example.  Its axiom audit uses only
\texttt{propext}, \texttt{Classical.choice}, and \texttt{Quot.sound}.

The algebraic part of S-I4C is therefore closed: a chain generated by actual
positive five-term moves has the required finite exterior boundary.  The
route remains open at the harder existence and analytic stages.  For an
arbitrary target cell, one must still construct a finite positive generated
filling satisfying useful bounds on chain length, coefficients, prime
support, and the two costs on the right of
\eqref{eq:sftbb-generated-cost}.  More precisely, put
\[
 \mathcal Q(\Gamma)=\sum_j|n_j|(M(P_j)+V(P_j)),
 \qquad
 \mathcal R(\Gamma)=\sum_j|n_j|R(P_j).
\]
For every $\epsilon>0$, Gate VF still requires constants
$K_\epsilon,L_\epsilon$ and a generated filling for each primitive target
of height $H=\log c$ and radical logarithm
$\rho=\log\operatorname{rad}(abc)$ such that
\[
 \mathcal R(\Gamma)\leq\epsilon\mathcal Q(\Gamma)+K_\epsilon H,
 \qquad
 \mathcal Q(\Gamma)\leq2H\rho+L_\epsilon H.
\]
Neither estimate nor such a uniform construction is proved here, and this
section proves neither the $abc$ conjecture nor its negation.
